\documentclass[11pt]{article}

\usepackage[utf8]{inputenc}
\usepackage[T1]{fontenc}
\usepackage{geometry}
\usepackage{amsmath}
\usepackage{amssymb}
\usepackage{booktabs}
\usepackage{hyperref}
\usepackage{xurl}
\usepackage{xcolor}
\usepackage{listings}
\usepackage{enumitem}
\usepackage{parskip}
\usepackage{graphicx}
\graphicspath{{../images/}}

\geometry{margin=1in}

\lstset{
  language=Java,
  basicstyle=\ttfamily\small,
  keywordstyle=\color{blue},
  commentstyle=\color{gray},
  stringstyle=\color{teal},
  showstringspaces=false,
  columns=fullflexible,
  frame=single,
  framerule=0.2pt,
  breaklines=true
}

\setlist[itemize]{topsep=0.3em,itemsep=0.2em}
\setlist[enumerate]{topsep=0.3em,itemsep=0.2em}

% Simple per-section source line.
\newcommand{\notesources}[1]{\par\noindent{\small\color{gray}\textit{Sources:} \sloppy #1\par}}

% Unit-wise source strings for this subject (used by slides/notes).
% Path is relative to the lecture `latex/` folder (the compilation CWD).
% OOPS with Java (CSEG1044) — unit-wise sources
%
% These are short, student-facing reference strings intended to be used with:
%   - \slidesources{...}  (in beamer slides)
%   - \notesources{...}   (in lecture notes)
%
% Style inspiration: other subjects in My-Subjects/ use semicolon-separated sources
% and include an "accessed" date for web documentation.

\newcommand{\OOPSUnitISources}{Schildt, \textit{Java: A Beginner's Guide} (10e), McGraw-Hill (Java fundamentals + OOP basics).; Oracle Java Tutorials: Getting Started + Language Basics + Classes and Objects (accessed \today).; Java SE API docs (e.g., \texttt{String}, \texttt{Scanner}) (accessed \today).}

\newcommand{\OOPSUnitIISources}{Schildt, \textit{Java: The Complete Reference} (12e), McGraw-Hill (classes, inheritance, packages, interfaces).; Bloch, \textit{Effective Java} (3e), Addison-Wesley, 2018 (design choices; interfaces vs abstract classes).; Oracle Java Tutorials: Classes and Objects + Interfaces + Packages (accessed \today).; Java SE API docs (e.g., \texttt{Comparable}, \texttt{Runnable}) (accessed \today).}

\newcommand{\OOPSUnitIIISources}{Schildt, \textit{Java: The Complete Reference} (12e) (nested/inner classes, exceptions, threads, I/O).; Oracle Java Tutorials: Nested Classes + Exceptions + Concurrency + Basic I/O (accessed \today).; Java SE API docs (e.g., \texttt{Thread}, \texttt{AutoCloseable}) (accessed \today).}

\newcommand{\OOPSUnitIVSources}{Schildt, \textit{Java: The Complete Reference} (12e) (generics, lambdas, Swing, JDBC).; Bloch, \textit{Effective Java} (3e) (generics + lambdas).; Oracle Java Tutorials: Generics + Lambda Expressions + Swing + JDBC Basics (accessed \today).; Java SE API docs (e.g., \texttt{java.util.function}, \texttt{javax.swing}, \texttt{java.sql}) (accessed \today).}

\newcommand{\OOPSUnitVSources}{Schildt, \textit{Java: The Complete Reference} (12e) (Collections Framework + wrapper classes).; Oracle Java docs: Collections Framework overview (accessed \today).; Java SE API docs (\texttt{java.util} collections; wrapper classes in \texttt{java.lang}) (accessed \today).}

\newcommand{\OOPSUnitVISources}{Git SCM: \textit{Pro Git} (2e) (version control workflow), accessed \today.; GitHub Docs (repositories, issues, pull requests), accessed \today.; Oracle Java Tutorials: Swing + JDBC (accessed \today).; JUnit 5 User Guide (accessed \today).; Apache Maven Project: Getting Started (accessed \today).; Gradle User Manual: Getting Started (accessed \today).; UML class diagrams: UML 2.x overview/spec (accessed \today).}

\newcommand{\OOPSMiscWhyInterfacesSources}{Schildt, \textit{Java: The Complete Reference} (12e) (interfaces).; Bloch, \textit{Effective Java} (3e) (prefer interfaces where appropriate).; Oracle Java Tutorials: Interfaces (accessed \today).; Java SE API docs (e.g., \texttt{Comparable}, \texttt{Runnable}) (accessed \today).}



\title{Object Oriented Programming (CSEG1044)\\Unit 05 -- Lecture 01 Notes\\Collections Framework Overview + Iteration + List}
\author{Tofik Ali}
\date{\today}

\begin{document}
\maketitle
\tableofcontents

\section{Lecture Overview}
The Collections Framework gives you standard, reusable data structures in Java.
Instead of creating arrays manually and writing custom logic, you use:
\begin{itemize}
  \item \texttt{List} for ordered sequences,
  \item \texttt{Set} for unique values,
  \item \texttt{Map} for key-value lookups.
\end{itemize}
This lecture focuses on \textbf{lists}, \textbf{iteration}, and \textbf{sorting}.
\notesources{\OOPSUnitVSources}

\section{Syllabus Alignment (Unit 05)}
\textbf{Unit 05:} Collections, iteration, Collection interface; List; sorting

\section{Learning Outcomes}
After this lecture, you should be able to:
\begin{itemize}
  \item Explain \texttt{Iterable}, \texttt{Collection}, and \texttt{Iterator}.
  \item Choose \texttt{ArrayList} vs \texttt{LinkedList} for common use cases.
  \item Sort a list using \texttt{Comparable} and \texttt{Comparator} (including lambdas).
\end{itemize}

\section{Key Terms (Quick)}
\begin{itemize}
  \item \textbf{Collection}: group of elements (root interface for many collections).
  \item \textbf{Iterable}: can be used in for-each loop (provides \texttt{iterator()}).
  \item \textbf{Iterator}: cursor used to traverse a collection safely.
  \item \textbf{List}: ordered collection with duplicates allowed and index-based access.
  \item \textbf{Comparable}: natural ordering (\texttt{compareTo}).
  \item \textbf{Comparator}: custom ordering (\texttt{compare}).
\end{itemize}

\section{Detailed Notes}
\subsection{Core interfaces: \texttt{Iterable}, \texttt{Collection}, \texttt{Iterator}}
\textbf{Iterable:}
\begin{itemize}
  \item Any \texttt{Iterable} can be used with the enhanced for loop.
\end{itemize}
\textbf{Iterator:}
\begin{itemize}
  \item Provides \texttt{hasNext()}, \texttt{next()}, and optional \texttt{remove()}.
  \item Use iterator when you need to remove elements while looping.
\end{itemize}

\paragraph{Common pitfall.} Removing from a list inside a for-each loop can throw \texttt{ConcurrentModificationException}. Use an \texttt{Iterator} instead.

\subsection{List: ArrayList vs LinkedList}
\textbf{ArrayList:}
\begin{itemize}
  \item Backed by a dynamic array.
  \item Fast random access: \texttt{get(i)} is efficient.
  \item Good for most use cases: iterate a lot, access by index.
\end{itemize}

\textbf{LinkedList:}
\begin{itemize}
  \item Nodes with links (doubly linked list).
  \item Slower random access (must traverse nodes).
  \item Useful when you frequently add/remove at ends and use it as a queue/deque.
\end{itemize}

\paragraph{Rule of thumb.} Use \texttt{ArrayList} by default unless you have a clear reason to use \texttt{LinkedList}.

\subsection{Sorting: Comparable vs Comparator}
\textbf{Comparable:}
\begin{itemize}
  \item implemented by the class itself (natural ordering).
  \item example: sort students by roll number.
\end{itemize}

\textbf{Comparator:}
\begin{itemize}
  \item separate ordering logic, reusable and flexible.
  \item example: sort students by CGPA descending.
\end{itemize}
\notesources{\OOPSUnitVSources}

\section{Worked Example: list + iterator + sorting}
\begin{lstlisting}
import java.util.*;

class Student {
    private final String rollNo;
    private final double cgpa;

    Student(String rollNo, double cgpa) {
        this.rollNo = rollNo;
        this.cgpa = cgpa;
    }

    public String getRollNo() { return rollNo; }
    public double getCgpa() { return cgpa; }
    public String toString() { return rollNo + ":" + cgpa; }
}

public class Demo {
    public static void main(String[] args) {
        List<Student> list = new ArrayList<>();
        list.add(new Student("21CS1002", 9.1));
        list.add(new Student("21CS1001", 8.2));
        list.add(new Student("21CS1003", 3.9));

        // Remove weak students safely using Iterator
        Iterator<Student> it = list.iterator();
        while (it.hasNext()) {
            if (it.next().getCgpa() < 4.0) it.remove();
        }

        // Sort by CGPA descending using Comparator + lambda
        list.sort((a, b) -> Double.compare(b.getCgpa(), a.getCgpa()));
        System.out.println(list);
    }
}
\end{lstlisting}

\section{In-Class Exercises (with brief solutions)}
\begin{enumerate}
  \item Why do we use \texttt{Iterator.remove()} instead of \texttt{list.remove(...)} inside a for-each loop?\par
  \textbf{Answer:} to avoid concurrent modification issues; iterator removal is safe during traversal.

  \item Which list is better for random access by index?\par
  \textbf{Answer:} \texttt{ArrayList}.

  \item What is the difference between \texttt{Comparable} and \texttt{Comparator}?\par
  \textbf{Answer:} \texttt{Comparable} defines natural order inside class; \texttt{Comparator} defines external/custom order.
\end{enumerate}

\section{Self-Study Practice (no solutions)}
\begin{enumerate}
  \item Create a list of integers and remove all odd numbers using an iterator.
  \item Sort a list of strings by length using a comparator lambda.
  \item Implement \texttt{Comparable} in a class \texttt{Book} to sort by title.
\end{enumerate}

\section{Checkpoint Questions (with sample answers)}
\textbf{Q1.} Why does Java provide a Collections Framework?\par
\textbf{Sample answer:} It provides standard, reusable data structures and algorithms, reducing bugs and improving productivity compared to writing custom structures every time.

\vspace{0.6em}
\textbf{Q2.} When should you prefer \texttt{Comparator} over \texttt{Comparable}?\par
\textbf{Sample answer:} Prefer \texttt{Comparator} when you need multiple sorting orders or when you cannot/should not modify the class to define natural ordering.

\end{document}
